%%%%%%%%%%%%%%%%%%%%%%%%%%%%%%%%%%%%%%%%%%%%%%%%%%%%%%%%%%%%%%%%%%%%%%%%%%%
%  SHT v8 — A Three-Layer Framework for Hallucination Detection in LLMs
%  Token Theorem + Claim Conformal Calibration + Empirical Hallucination Bridge
%
%  Target venue: NLPCC 2026 (Springer LNCS format)
%  Format:       LLNCS 2.21, 12-page limit (incl. references)
%  Status:       SUBMISSION — All numerical results from Layer 3 runs injected
%
%  ANONYMIZATION NOTE for double-blind review:
%    - Remove \inst{} affiliation details, acknowledgments
%    - Replace author names with "Anonymous"
%    - Remove or mask pipeline artifact URLs if any
%%%%%%%%%%%%%%%%%%%%%%%%%%%%%%%%%%%%%%%%%%%%%%%%%%%%%%%%%%%%%%%%%%%%%%%%%%%
\documentclass[runningheads]{llncs}

% ── Fonts ────────────────────────────────────────────────────────────────
\usepackage[T1]{fontenc}

% ── Math ──────────────────────────────────────────────────────────────────
\usepackage{amsmath,amssymb}
\usepackage{bm}

% ── Figures / Color ───────────────────────────────────────────────────────
\usepackage{graphicx}
\usepackage{xcolor}

% (TODO marker command removed for submission — all results injected)

% ═════════════════════════════════════════════════════════════════
%  Math notation shortcuts
% ═════════════════════════════════════════════════════════════════
\newcommand{\qt}{q_t}
\newcommand{\pt}{p_t}
\newcommand{\qxt}{q_t(X_t)}
\newcommand{\qlsq}{\|q_t\|_2^2}
\newcommand{\stscore}{s_t}
\newcommand{\Mt}{M_t}
\newcommand{\Lambdacohort}{\Lambda_t^{(j)}}
\newcommand{\claimscore}{A_{\text{claim}}}
\newcommand{\claimscorelog}{A_{\text{claim}}^{\log}}
\newcommand{\Hnull}{H_0^{\text{stat}}}
\newcommand{\Talpha}{T_\alpha}
\newcommand{\betacohort}{\beta_t^{(j)}}

%%%%%%%%%%%%%%%%%%%%%%%%%%%%%%%%%%%%%%%%%%%%%%%%%%%%%%%%%%%%%%%%%%%%%%%%%%%
\begin{document}

% ── Title ─────────────────────────────────────────────────────────────
\title{SHT: A Self-Surprise Conformal Framework for LLM Hallucination Analysis\thanks{SHT is the framework identifier retained from previous versions of this work.}}
\titlerunning{SHT: Self-Surprise Conformal Framework}

% ── Author (anonymize for double-blind review) ────────────────────────
\author{Jie Yuan\inst{1}}

\authorrunning{J. Yuan}

\institute{Kean University, Union NJ 07083, USA \\
\email{yuanjie@kean.edu}}

\maketitle

% ── Abstract ──────────────────────────────────────────────────────────
\begin{abstract}
We propose SHT, a conformal-calibrated framework for analyzing
hallucinations in large language models (LLMs) via token-derived
self-surprise statistics. The framework has three components.
We define a centered-Brier self-surprise residual
$s_t = \|q_t\|_2^2 - q_t(X_t)$ that is bounded in $[-1, 1]$ and
mean-zero under the model's own (post-temperature, post-top-$p$)
sampling distribution $q_t$---the trivial statistical null
$\Hnull: X_t \mid \mathcal{F}_{t-1} \sim q_t$. Aggregating $s_t$ over
a claim's token span yields a per-claim score; we then form a per-trace
summary score and calibrate via per-category split conformal prediction at
the trace level for a finite-sample bound
$\mathbb{P}(\bar{A}_{\text{new}} > \bar{T}_{\alpha,c} \mid Y_{\text{new}} =
\text{factual}, c) \le \alpha + 1/(n_c^{\text{tr}}+1)$ on the false alarm
rate over factual traces, conditional on category, under trace-level
within-category exchangeability (Theorem~\ref{thm:layer2}, where
$n_c^{\text{tr}}$ is the number of factual calibration traces in category
$c$). The trace-level unit matches our cal/test split protocol exactly
and delivers an exact finite-sample guarantee. A claim-level analogue
threshold using $\alpha + 1/(m_c+1)$ as an informal reference value
(Remark~\ref{prop:claim-level}) is used for individual-claim AUROC
characterization, with empirical Clopper--Pearson per-prompt miscoverage
diagnostics flagging $(c, \text{prompt})$ cells whose lower CI exceeds
the informal reference. We then characterize, per
category, the empirical relationship between this score and
claim-level hallucination labels using AUROC, PR-AUC,
Kolmogorov--Smirnov, and Mann--Whitney $U$ with trace-level bootstrap
confidence intervals. The comparison includes mean-entropy, mean-NLL,
max-NLL, and per-claim variance baselines, so that any advantage
attributable to the proposed score is isolated from baseline
second-moment readouts. As a structural byproduct, the cumulative
product $M_t = \prod_t(1 + \beta\, s_t)$ is an anytime-valid e-process
under $\Hnull$ (Theorem~\ref{thm:layer1}). We characterize the
behavior of $M_t$ under repeated peeking in two streaming experiments
on identical input streams: (a) a delay-recall comparison against a
CUSUM baseline run on the same score $s_t$, characterizing the
contribution of multiplicative-vs-additive aggregation; and (b) a
three-arm peeking-FAR experiment on the same $M_t$, contrasting the
Ville-derived $1/\alpha$, a tight empirically-calibrated anytime-valid
threshold $c_A^*$ (the per-category $(1-\alpha)$-empirical-quantile of
$\sup_t M_t$ over cal-factual traces in each category $c$), and a
single-stop critical value $c^*$. By the per-category structural ordering
$c^* \leq c_A^* \leq 1/\alpha$ (see Remark~\ref{rem:theorem1-scope}),
the resulting FAR gaps are reported as
joint characterizations of (anytime-valid status, calibration-target
choice) rather than as isolated attribution to anytime-validity; both
experiments and the batch comparison are reported in the main text
regardless of outcome.
We distinguish Channel~U (uncertainty-driven hallucinations,
structurally visible to token-level proper scoring rules) from
Channel~O (overconfident hallucinations, structurally invisible) as a
methodological decomposition of detectability boundaries; Channel~O
is declared a documented limitation rather than a hidden one. Three
further limitations are disclosed explicitly: (i) ground-truth labels
are derived from three LLM annotators, so per-category AUROC measures
agreement between LLM-derived signals on a unanimous-consensus subset
rather than verification against a human gold standard; (ii) the
five-category design (C1--C5) yields empirically interpretable
results predominantly for C5, with C1--C4 reported as supporting
characterization under acknowledged power and structural limitations;
(iii) the unanimity merge rule retains the cases on which all three
annotators reach consensus, which systematically overestimates
detection performance on the underlying claim population.
Experiments on Qwen2.5-7B-Instruct with 25~prompts $\times$
5~seeds~$=$ 125~generation trajectories, 1{,}026~atomic claims, and
three independent LLM annotators (DeepSeek V4 Pro, GLM-4.7, Claude
Sonnet 4.6) under a unanimity merge rule yield
a labeled dataset of
667~3-annotator-unanimous claims (560 factual, 107 hallucinated; Fleiss'
$\kappa = 0.5763$).
On batch evaluation, the e-process is comparable to NLL baselines
($\Delta_{\text{AUROC}} = -0.103$, 95\% CI $[-0.246, +0.101]$, all four
non-degenerate length quartiles negative).
Streaming delay-recall is not evaluated at $n_h = 8$ test-hallucinated
trace-category records; the evaluated hallucinated subset is Channel-O
leaning (see \S\ref{sec:channels}), placing it outside the framework's
detection scope by construction.
Joint-characterization peeking-FAR analysis is reported with finite-sample
caveats due to per-category calibration counts
$n_c^{\text{tr}} \in \{13,14,17,17,17\}$ below the non-vacuous-quantile
floor $\lceil 1/\alpha \rceil + 1 = 21$ at $\alpha = 0.05$; empirical
$\hat{\text{FAR}}_A = 0/31$ does not violate the Ville bound but cannot
certify it tighter than $[0, 0.112]$ at this sample size.
\keywords{Hallucination analysis \and conformal prediction \and
sequential testing \and self-surprise \and anytime-valid inference
\and large language models.}
\end{abstract}

% ═══════════════════════════════════════════════════════════════════════
\section{Introduction}

\subsection{The Problem}

Large language models generate fluent text but are prone to
\textit{hallucinations}---statements that are factually incorrect,
unsubstantiated, or internally inconsistent. Existing detection methods
fall into three broad classes:
\begin{enumerate}
    \item \textbf{Multi-sample consistency} methods (e.g., SelfCheckGPT~\cite{manakul2023selfcheck})
          detect factuality degradation via repeated sampling, which is
          computationally expensive and post-hoc;
    \item \textbf{Semantic-entropy / confabulation} methods
          (e.g.,~\cite{farquhar2024detecting}) identify confabulation subsets
          but require claim-level modeling;
    \item \textbf{External knowledge grounding} (e.g., FActScore~\cite{min2023factscore})
          verifies claims against retrieved evidence, introducing retrieval
          latency and coverage gaps.
\end{enumerate}
None of these methods operate \textit{online} at the token level with formal
statistical guarantees on false alarm rates.

\subsection{Core Insight: Reframing the Null Hypothesis}

A central conceptual shift underlies our framework. We do \textbf{not} use the
e-process to ``reject the hypothesis that the model is truthful.'' Instead, we
define the null hypothesis purely in terms of the model's own sampling behavior:
\begin{equation}\label{eq:hnull}
\Hnull:\; X_t \mid \mathcal{F}_{t-1} \sim q_t,
\end{equation}
where $q_t$ is the \textit{actual} sampling distribution after temperature
scaling and top-$p$ truncation. Under standard HuggingFace sampling,
$X_t \sim q_t$ is a \textit{behavioral fact}, not a falsifiable statistical
assumption. The e-process therefore functions as a \textbf{rare-trajectory
detector}: it flags token paths that are improbable under the model's own
sampling regime, without making any claim about factual truth.

This reframing has an important consequence: the e-process can flag
``self-surprise'' even when the model is generating factually correct text.
Conversely, the e-process can remain silent when the model confidently emits a
falsehood (Channel~O hallucinations). The connection between self-surprise and
factual error is therefore \textbf{empirical}, not logical---and quantifying
this connection is the role of Layer~3.

\subsection{The Three-Layer Framework}

Our framework consists of three structurally independent layers with
non-overlapping mathematical foundations:
\begin{itemize}
    \item \textbf{Layer~1 (Token Theorem)}: An anytime-valid e-process over
          token-level self-surprise scores, with Type~I error control under
          $\Hnull$ (Theorem~\ref{thm:layer1}).
    \item \textbf{Layer~2 (Claim Conformal Calibration)}: Split conformal
          prediction over atomic claims to bound the false alarm rate on
          factual claims, relying on exchangeability rather than $\Hnull$
          (Theorem~\ref{thm:layer2}).
    \item \textbf{Layer~3 (Empirical Hallucination Bridge)}: An empirical
          study connecting token-level self-surprise signals to claim-level
          hallucination labels, evaluated per category with bootstrap
          confidence intervals.
\end{itemize}

\subsection{Two Hallucination Channels}

We explicitly categorize hallucinations into two channels:
\begin{itemize}
    \item \textbf{Channel~U (Uncertainty-driven)}: The model samples in a
          low-knowledge region, producing tokens with
          $q_t(X_t) < \|q_t\|_2^2$. The e-process is structurally sensitive
          to this category.
    \item \textbf{Channel~O (Overconfident hallucination)}: The model
          confidently outputs a false token with high $q_t(X_t)$---often
          due to a false premise embedded in the prompt or context. Token-level
          proper scoring rules are structurally blind to this category.
          Addressing Channel~O requires orthogonal methods and is declared
          as explicit future work.
\end{itemize}

\subsection{Contributions}
\begin{enumerate}
    \item A token-derived self-surprise statistic
          $s_t = \|q_t\|_2^2 - q_t(X_t)$ with explicit boundedness and
          mean-zero properties under $\Hnull$, suitable for both batch
          claim-level scoring and sequential streaming evaluation
          (Lemma~\ref{lem:mean-zero}, Lemma~\ref{lem:bounded}).
    \item A trace-level per-category split conformal calibration
          procedure (Theorem~\ref{thm:layer2}) that delivers an exact
          finite-sample false-alarm-rate bound on factual traces
          conditional on category, where the unit of exchangeability
          (trace) matches the unit of the cal/test split. A claim-level
          analogue threshold $T_{\alpha,c}$ (Remark~\ref{prop:claim-level})
          is reported as an \emph{informal reference} using
          $\alpha + 1/(m_c+1)$, with per-prompt Clopper--Pearson
          miscoverage cells flagged as a diagnostic rather than as a
          formal bound (Remark~\ref{rem:trace-unit-rationale}).
    \item An empirical batch characterization on 1{,}026~atomic claims
          across five categories, comparing the proposed score against
          mean-entropy, mean-NLL, max-NLL, and \emph{per-claim variance}
          baselines, with a length-stratified sanity check; honest
          §6 limitations disclose that the effective empirical
          contribution is dominated by C5 and that the underlying
          labels are LLM-derived.
    \item A streaming delay-recall comparison against a CUSUM baseline
          run on the same score $s_t$, characterizing the contribution
          of the multiplicative aggregation mechanism relative to
          additive cumulation on the same underlying signal.
    \item An empirical characterization of false alarm rate behavior
          under repeated peeking across three threshold regimes on
          factual traces, providing operational context for
          Theorem~\ref{thm:layer1}'s structural guarantee: the same
          single-cohort wealth process $M_t = \prod_t(1 + \beta s_t)$
          is monitored under (i)~the Ville-derived $1/\alpha$
          (conservative anytime-valid threshold), (ii)~a tight
          empirically-calibrated anytime-valid threshold $c_A^*$ equal
          to the per-category $(1-\alpha)$-empirical-quantile of
          $\sup_t M_t$ over cal-factual traces in each category $c$,
          and (iii)~a single-stop critical value $c^*$ calibrated to
          $\mathbb{P}(M_{T_i}\!\ge\!c^*)\!=\!\alpha$ at each trace's
          natural terminal. Denoting the three alarms by Arm~A,
          Arm~$A'$, and Arm~B for thresholds $1/\alpha$, $c_A^*$, and
          $c^*$ respectively (precise definitions in
          Remark~\ref{rem:theorem1-scope}), by the per-category
          structural ordering $c^* \leq c_A^* \leq 1/\alpha$ (where
          $c^* \leq c_A^*$ follows from $\sup_t M_t \geq M_{T_i}$
          pointwise dominance on the same category-$c$ cal-factual
          sample, and $c_A^* \leq 1/\alpha$ follows from Ville's
          inequality under $\Hnull$;
          see Remark~\ref{rem:theorem1-scope}), the threshold
          magnitudes are jointly determined by both the
          calibration-target choice (sup-over-trajectory vs.\
          terminal-only) and the anytime-valid property; we therefore
          report the resulting FAR gaps as joint characterizations
          rather than as isolated attribution to any single property.
          By construction, whenever the peek set $K$ includes each
          trace's natural terminal time $T_i$ (the convention we adopt
          throughout), the lower bound $\text{FAR}_B(K) \geq \alpha$
          holds (Arm B's threshold $c^*$ is calibrated to alarm at rate
          $\alpha$ at $T_i$). Monotonicity
          $K' \supseteq K \Rightarrow \text{FAR}_B(K') \geq
          \text{FAR}_B(K)$ holds by set inclusion without
          distributional assumption.
    \item A demarcation of Channel~U vs.\ Channel~O as a methodological
          decomposition of detectability boundaries, with Channel~O
          declared a documented structural limitation.
\end{enumerate}

% ═══════════════════════════════════════════════════════════════════════
\section{Preliminaries}

\subsection{Strictly Proper Scoring Rules and Bregman Divergence}

A scoring rule $S(q, x)$ maps a predicted distribution $q$ over a finite
vocabulary $\mathcal{V}$ and an outcome $x \in \mathcal{V}$ to a real-valued
loss. $S$ is \textit{strictly proper} if the expected score
$\bar{S}(q, p) := \mathbb{E}_{Y \sim p}[S(q, Y)]$ satisfies
$\bar{S}(q, p) \ge \bar{S}(p, p)$ with equality iff $q = p$. Every strictly
proper scoring rule induces a Bregman divergence
$D_\phi(p \| q) = \bar{S}(q, p) - \bar{S}(p, p)$, where
$\phi(p) := -\bar{S}(p, p)$ is a strictly convex generator
function~\cite{gneiting2007strictly}.

We use the \textbf{Brier score}:
\begin{equation}
\text{BS}(q, x) = \sum_{v \in \mathcal{V}} (\mathbf{1}_{\{v=x\}} - q(v))^2
                = 1 - 2q(x) + \|q\|_2^2,
\end{equation}
with generator $\phi(p) = \|p\|_2^2$ and Bregman divergence $\|p - q\|_2^2$.

\subsection{E-Processes and Ville's Inequality}

\begin{definition}
A nonnegative process $(E_t)_{t \ge 0}$ adapted to a filtration
$(\mathcal{F}_t)$ is an e-process for a family of distributions $\mathcal{P}$
if $E_0 \le 1$ and for every $P \in \mathcal{P}$ and every
$(\mathcal{F}_t)$-stopping time $\tau$,
$\mathbb{E}_P[E_\tau] \le 1$.
\end{definition}

\begin{theorem}[Ville's inequality~\cite{ville1939etude}]
If $(E_t)$ is a nonnegative supermartingale with $E_0 \le 1$, then for any
$c > 0$, $\mathbb{P}(\sup_t E_t \ge c) \le E_0 / c \le 1/c$. In particular,
for $\tau = \inf\{t: E_t \ge 1/\alpha\}$,
$\mathbb{P}(\tau < \infty) \le \alpha$.
\end{theorem}

\subsection{Cohort E-Process and E-Detector}

Following~\cite{shin2023edetectors}, for each start time $j \ge 1$ we define
a \textit{cohort e-process} $\Lambda_t^{(j)}$ that begins accumulating evidence
at time $j$. The mixture e-process
$M_t = \sum_{j} w_j \Lambda_t^{(j)}$ weights cohorts by deterministic
nonnegative weights $\{w_j\}$ with $\sum_j w_j \le 1$.
An important practical detail: inactive cohorts (those with $j > t$) contribute
$w_j \cdot 1$ to the mixture, so
$M_t = \sum_{j \le t} w_j \Lambda_t^{(j)} + \sum_{j > t} w_j$. The tail sum
must be preserved for the supermartingale property to hold.

\subsection{Split Conformal Prediction}

Given a calibration set $\{A_1, \ldots, A_m\}$ of nonconformity scores for
factual claims, and a new exchangeable factual claim with score
$A_{\text{new}}$, the split conformal
framework~\cite{vovk2005algorithmic,lei2018distribution} guarantees:
\begin{equation}
\mathbb{P}\!\left(A_{\text{new}} > T_\alpha \right) \le \alpha + \frac{1}{m+1},
\end{equation}
where $T_\alpha$ is the
$\lceil (m+1)(1-\alpha) \rceil$-th order statistic of
$\{A_1, \ldots, A_m\}$.

% ═══════════════════════════════════════════════════════════════════════
\section{Setup and Assumptions}

\subsection{Generation Process and Filtration}

Let $\mathcal{V}$ be the finite vocabulary. At each step $t$, the LLM
computes logits $\ell_t$ and raw softmax $p_t$ from the context history.
The \textit{actual sampling distribution} $q_t$ is obtained from $p_t$ by
applying temperature $T$ and top-$p$ truncation:
$q_t(v) \propto p_t(v)^{1/T}$ restricted to the minimal set of tokens with
cumulative probability $\ge p$. The next token $X_t$ is sampled
$X_t \sim q_t$. The natural filtration is
$\mathcal{F}_t = \sigma(X_1, \ldots, X_t)$, and both $p_t$ and $q_t$ are
$\mathcal{F}_{t-1}$-measurable.

\begin{remark}[Implementation-critical]
All downstream theory uses $q_t$, not $p_t$. If the generation pipeline only
saves $p_t$ under $T \neq 1$ or $\texttt{top-p} < 1$, the centered token-level
score is no longer mean-zero, and the e-process property is lost.
Our pipeline explicitly saves $q_t(X_t)$ and $\|q_t\|_2^2$ for every token.
\end{remark}

\subsection{Two Distinct Null Hypotheses}

\begin{definition}[Statistical null]
$\Hnull: X_t \mid \mathcal{F}_{t-1} \sim q_t$ for all $t$.
\end{definition}

\begin{definition}[Factual null]
$H_0^{\text{fact}}$: a claim spanning tokens $[a, b]$ is factually correct.
\end{definition}

\begin{remark}
Under standard HuggingFace sampling, $X_t \sim q_t$ is a behavioral fact, not
a falsifiable assumption. Even when the generated claims are entirely
hallucinated, $\Hnull$ continues to hold as long as the sampling mechanism
follows $q_t$. Therefore, an e-process based on $\Hnull$ cannot inflate in
expectation merely because the model hallucinates---it can only detect that
the \textit{realized trajectory} lies in a low-density region of $q_t$.
\end{remark}

\begin{remark}[OOD prompts]
Out-of-distribution prompts change the shape of $q_t$ (shifted support,
concentrated mass, etc.) but do not change the fact that $X_t \sim q_t$.
OOD effects are therefore \textit{empirical}---they may make hallucination
and low-support sampling statistically correlated, which is a Layer~3
phenomenon, not a Layer~1 violation of $\Hnull$.
\end{remark}

\subsection{Channel~U and Channel~O}
\label{sec:channels}

\begin{itemize}
    \item \textbf{Channel~U (Uncertainty-driven)}: The model operates in an
          uncertain region, producing tokens where
          $q_t(X_t) < \|q_t\|_2^2$ is sustained. The e-process is structurally
          designed to detect this channel.
    \item \textbf{Channel~O (Overconfident hallucination)}: The model
          confidently outputs a false token with high $q_t(X_t)$ and
          $s_t \le 0$, so the e-process does not trigger. Causes include
          false premises embedded in the prompt, memorized falsehoods, or
          pattern-completion artifacts. Channel~O requires orthogonal methods
          and is explicit future work.
\end{itemize}

\paragraph{Dataset characterization.}
On the present dataset, hallucinated atomic claims distribute by category
as $C_1{=}7$, $C_2{=}2$, $C_3{=}41$, $C_4{=}27$, $C_5{=}30$ (total~107
hallucinated of 667 unanimous-merged claims). The post-cutoff-knowledge
category $C_5$ alone accounts for 28\% of hallucinations, and the wider
Channel-O proxy $C_1 + C_5$ accounts for 34.6\% — placing this dataset
in a Channel-O-leaning regime per the above taxonomy.

% ═══════════════════════════════════════════════════════════════════════
\section{Layer~1: Token-Level Self-Surprise E-Process}

\subsection{Score Function: Centered Brier Residual}

We define the token-level score as the centered Brier residual:
\begin{equation}\label{eq:st}
s_t = \tfrac{1}{2}\bigl(\text{BS}(q_t, X_t) -
      \mathbb{E}_{Y \sim q_t}[\text{BS}(q_t, Y)]\bigr)
    = \|q_t\|_2^2 - q_t(X_t) \in [-1, 1].
\end{equation}

\begin{lemma}[Mean-zero under $\Hnull$]\label{lem:mean-zero}
Under $\Hnull$, $\mathbb{E}[s_t \mid \mathcal{F}_{t-1}] = 0$.
\end{lemma}

\begin{proof}
Since $X_t \mid \mathcal{F}_{t-1} \sim q_t$ and $q_t$ is
$\mathcal{F}_{t-1}$-measurable,
$\mathbb{E}[q_t(X_t) \mid \mathcal{F}_{t-1}]
 = \sum_v q_t(v) \cdot q_t(v) = \|q_t\|_2^2$.
Therefore $\mathbb{E}[s_t \mid \mathcal{F}_{t-1}]
 = \|q_t\|_2^2 - \|q_t\|_2^2 = 0$.
\end{proof}

\begin{lemma}[Boundedness]\label{lem:bounded}
$s_t \in [-1, 1]$ almost surely.
\end{lemma}

\begin{proof}
$q_t$ is a probability distribution over the finite vocabulary
$\mathcal{V}$, so $q_t(v) \in [0, 1]$ for every $v$ and
$\|q_t\|_2^2 = \sum_v q_t(v)^2 \in [0, 1]$. Therefore
$s_t = \|q_t\|_2^2 - q_t(X_t) \in [-1, 1]$ almost surely.
\end{proof}

\subsection{Why the Expectation Is Not a Bregman Divergence}

A correction from a prior version (v7): under an arbitrary alternative
distribution $X_t \sim r_t \neq q_t$,
\begin{equation}\label{eq:expectation-r}
\mathbb{E}_{X_t \sim r_t}[s_t \mid \mathcal{F}_{t-1}]
    = \|q_t\|_2^2 - \langle r_t, q_t \rangle
    = \langle q_t - r_t,\, q_t \rangle,
\end{equation}
which is \textbf{not} $\|r_t - q_t\|_2^2$ and is not guaranteed to be
nonnegative. This is not a bug---it is the mathematical reason why
Channel~O hallucinations (where $r_t$ concentrates mass on the same tokens as
$q_t$) are invisible to the e-process. The Bregman divergence only appears in
the \textit{propriety gap} of the full Brier score, not in the centered
residual of~\eqref{eq:st}.

\subsection{Cohort E-Process Construction}

For each start index $j \ge 1$, define the cohort process:
\begin{equation}\label{eq:cohort}
\Lambda_t^{(j)} =
\begin{cases}
1, & t < j, \\[6pt]
\displaystyle\prod_{i=j}^{t} \bigl(1 + \beta_i^{(j)} s_i\bigr), & t \ge j,
\end{cases}
\end{equation}
where $\beta_i^{(j)} \in [0, B]$ with $B \in (0, 1)$, and crucially,
$\beta_i^{(j)}$ must be $\mathcal{F}_{i-1}$-measurable---it cannot depend on
$s_i$ or $X_i$.

\begin{lemma}[Cohort martingale]\label{lem:cohort-martingale}
Under $\Hnull$, each $\Lambda^{(j)}$ is a nonnegative martingale with
$\mathbb{E}[\Lambda_t^{(j)}] = 1$.
\end{lemma}

\begin{proof}
$1 + \beta_i^{(j)} s_i \ge 1 - B > 0$ ensures nonnegativity. For $t \ge j$,
$\mathbb{E}[\Lambda_t^{(j)} \mid \mathcal{F}_{t-1}]
 = \Lambda_{t-1}^{(j)}\bigl(1 + \beta_t^{(j)}\mathbb{E}[s_t \mid
   \mathcal{F}_{t-1}]\bigr)
 = \Lambda_{t-1}^{(j)}$ by Lemma~\ref{lem:mean-zero} and
$\beta_t^{(j)} \in \mathcal{F}_{t-1}$.
\end{proof}

\subsection{Mixture E-Process}

Let $\{w_j\}_{j \ge 1}$ be deterministic nonnegative weights with
$\sum_{j=1}^\infty w_j \le 1$. The mixture e-process is:
\begin{equation}\label{eq:mixture}
M_t = \sum_{j=1}^\infty w_j \Lambda_t^{(j)}
    = \underbrace{\sum_{j=1}^{t} w_j \Lambda_t^{(j)}}_{\text{active}}
    + \underbrace{\sum_{j=t+1}^{\infty} w_j}_{\text{inactive tail}}.
\end{equation}

\begin{remark}[Inactive tail]
In implementation, it is common to compute only the active sum, but then
$M_0^{\text{active}} = 0 \not\le 1$ and the Ville bound does not apply
directly. The correct procedure initializes $M \leftarrow \sum_j w_j$ and
moves weight from the inactive tail to active cohorts as new cohorts start.
\end{remark}

\subsection{Main Theorem}

\begin{theorem}[Anytime-valid Type~I control]\label{thm:layer1}
Let $M_t$ be defined by~\eqref{eq:mixture} with $\sum_j w_j \le 1$ and
$\beta_i^{(j)} \in [0, B] \subset [0, 1)$ being
$\mathcal{F}_{i-1}$-measurable. Define
$\tau_\alpha = \inf\{t: M_t \ge 1/\alpha\}$. Then
\[
\mathbb{P}_{\Hnull}\!\left(\tau_\alpha < \infty\right)
    = \mathbb{P}_{\Hnull}\!\left(\sup\nolimits_t M_t \ge 1/\alpha\right)
    \le \alpha.
\]
\end{theorem}

\begin{proof}
By Lemma~\ref{lem:cohort-martingale} and linearity of expectation,
$M_t$ is a nonnegative martingale with $M_0 = \sum_j w_j \le 1$.
Ville's inequality with $c = 1/\alpha$ gives the result.
\end{proof}

\begin{remark}[Scope of Theorem~\ref{thm:layer1}]
\label{rem:theorem1-scope}
Theorem~\ref{thm:layer1} controls
$\mathbb{P}_{\Hnull}(\sup_t M_t \ge 1/\alpha) \le \alpha$, where
$\Hnull: X_t \mid \mathcal{F}_{t-1} \sim q_t$ holds by construction
whenever the model samples from its own (post-temperature,
post-top-$p$) distribution. Two consequences follow.

\textbf{(i)~The guarantee is calibration, not detection power.}
The inequality bounds the Type~I rate of the rare-path monitor under
a null that is satisfied on \emph{every} trajectory the model
produces---factual or hallucinated. Theorem~\ref{thm:layer1}
therefore certifies that the monitor will not falsely alarm at rate
exceeding $\alpha$; it does \emph{not} imply that the monitor's alarm
rate on hallucinated trajectories is any higher than on factual ones.
In particular, Theorem~\ref{thm:layer1} does NOT claim
(a)~alarm~$\Rightarrow$~hallucination,
(b)~hallucination~$\Rightarrow$~$\Hnull$ violation, or
(c)~bounded false alarm rate on factual claims (this last is
addressed by Theorem~\ref{thm:layer2} under a separate
exchangeability assumption, not by Theorem~\ref{thm:layer1}).
Whether the score function $s_t$ empirically separates factual from
hallucinated trajectories is an empirical question, addressed in the
batch claim-level results and the streaming delay-recall results.

\textbf{(ii)~The behavior of $M_t$ under repeated peeking is
empirically characterized via a three-arm comparison on the same
wealth process.} The same single-cohort wealth process
$M_t = \prod_t(1 + \beta s_t)$ is monitored under three threshold
regimes that share the same input data and the same peek schedule:
the Ville-derived $1/\alpha$ (conservative anytime-valid threshold,
distribution-free), a tight empirically-calibrated anytime-valid
threshold $c_A^*$ equal to the $(1-\alpha)$-empirical-quantile of
$\sup_t M_t$ over cal-factual traces in category $c$ (which delivers,
by split conformal applied to the trace-level statistic $\sup_t M_t$
under cal/test trace-level exchangeability within category $c$, the
bound $\mathbb{P}_\text{test}(\sup_t M_t \ge c_A^* \mid c) \le \alpha
+ 1/(n_c^{\text{tr}}+1)$; subset containment $\{M_\tau \ge c_A^*\}
\subseteq \{\sup_t M_t \ge c_A^*\}$ for any stopping time $\tau$ then
extends this bound to any peek schedule), and a
single-stop critical value $c^*$ calibrated to
$\mathbb{P}(M_{T_i}\!\ge\!c^*)\!=\!\alpha$ at each trace's natural
terminal (\emph{not} anytime-valid under peeking). We denote by Arm~A,
Arm~$A'$, and Arm~B the alarms triggered by thresholds $1/\alpha$,
$c_A^*$, and $c^*$ respectively, and write
$\text{FAR}_X(K) := \mathbb{P}_{\text{factual, test}}\bigl(\sup_{\tau \in K} M_\tau \ge \text{threshold}_X\bigr)$
for $X \in \{A, A', B\}$ over a peek schedule $K$. By the per-category
structural ordering $c^* \leq c_A^* \leq 1/\alpha$, where
$c^* \leq c_A^*$ follows from $\sup_t M_t \geq M_{T_i}$ pointwise
dominance applied to the same category-$c$ cal-factual sample, and
$c_A^* \leq 1/\alpha$ follows from Ville's inequality applied to the
non-negative martingale $M_t$ under $\Hnull$ via the empirical
$(1-\alpha)$-quantile-to-tail correspondence on cal-factual traces, the
three threshold
magnitudes are jointly determined by both the calibration-target
choice (sup vs.\ terminal) and the anytime-valid property. We
therefore report the resulting gaps $\text{FAR}_B(K) -
\text{FAR}_{A'}(K)$ and $\text{FAR}_{A'}(K) - \text{FAR}_A(K)$ as
\emph{joint characterizations} rather than as isolated attribution to
any single property; the Ville bound's universal distribution-free
guarantee on $\text{FAR}_A$ remains exact at any peek schedule.

This separation between \emph{what the theorem certifies}
(calibration under a constructed null, robust to optional stopping)
and \emph{what the experiments establish} (empirical characterization
of FAR behavior under peeking across three threshold regimes, with
the Ville bound's exact anytime-valid Type-I control invariant under
peeking) is made explicit throughout the paper; the theorem is not
used as evidence for detection capability.
\end{remark}

\subsection{Betting Strategy}

The choice of $\beta_t^{(j)}$ is critical for power but must respect
$\mathcal{F}_{t-1}$-measurability. We consider:
\begin{itemize}
    \item \textbf{Fixed} $\beta = 0.1$: simple, no look-ahead, serves as baseline.
    \item \textbf{EMA plug-in}: exponential moving average of recent
          $|s_{t-k}|$ to adapt bet size while maintaining measurability.
    \item \textbf{ONS} (Online Newton Step)~\cite{hazan2007logarithmic}:
          adaptive betting with regret guarantees, applied only to past scores
          $s_1, \ldots, s_{t-1}$ to preserve $\mathcal{F}_{t-1}$-measurability;
          deferred from the main runs and reported only as an ablation per
          \texttt{phase\_d/PIPELINE\_MODIFICATIONS\_PLAN\_v2.md} §11.
\end{itemize}
Weight schedules: \textbf{harmonic}
$w_j = 1/(j(j+1))$ ($\sum_j w_j = 1$, $O(t^2)$) and
\textbf{dyadic}~\cite{shin2023edetectors}
($w_j > 0$ only for $j \in \{1, 2, 4, 8, \ldots\}$,
$\sum_j w_j = 1$, $O(t \log t)$).

The primary betting/weight combination for the main results is committed
to \texttt{phase\_d/runs/<ts>/manifest.json} before execution per the
pre-registration discipline in
\texttt{phase\_d/PIPELINE\_MODIFICATIONS\_PLAN\_v2.md} §6.4; all listed
variants are evaluated in the Group~D ablation
(\S\ref{ssec:results-groupD}).

% ═══════════════════════════════════════════════════════════════════════
\section{Layer~2: Claim-Level Conformal Calibration}

\subsection{Motivation}

Theorem~\ref{thm:layer1} controls
$\mathbb{P}_{\Hnull}(\text{alarm}) \le \alpha$, but this is not the same as
$\mathbb{P}(\text{alarm} \mid \text{claim is factual}) \le \alpha$---the two
probabilities live in different probability spaces. To bound the false alarm
rate on factual claims, we need an additional statistical principle:
split conformal calibration~\cite{vovk2005algorithmic,lei2018distribution}.

\subsection{Localized Claim Score}

A claim occupies token span $[a, b]$. We cannot use the global $M_t$ as the
claim score because $M_t$ accumulates from $t=1$, and a high $M_t$ during
$[a, b]$ may be inherited from evidence accumulated before $a$. We define two
localized aggregators:
\begin{align}
\claimscore([a, b]) &=
   \max_{j \in [a, b]} \max_{t \in [j, b]} w_j \Lambda_t^{(j)},
   \label{eq:aclaim} \\[6pt]
\claimscorelog([a, b]) &=
   \max_{t \in [a, b]} \bigl(\log M_t - \log M_{a-1}\bigr).
   \label{eq:aclaillog}
\end{align}
Equation~\eqref{eq:aclaim} is the localized e-value form (maximum over
cohorts starting within $[a, b]$). Equation~\eqref{eq:aclaillog} is the
log-wealth increment over the claim span---it is no longer an e-value but
empirically yields better ranking performance (higher AUROC). We report
both. Note that because $M_t$ is a sum (not product) over cohorts,
$\log M_t - \log M_{a-1}$ can \emph{underestimate} the evidence contributed
by cohorts starting within $[a, b]$ when pre-$a$ cohorts dominate the
mixture---the opposite failure mode from the multiplicative claim score
$\claimscore$, which can be inflated by short, high-$\Lambda$ cohorts
started mid-claim. We document these distinct pathologies in the ablation
report (Group~D).

\subsection{Per-Category Conformal Threshold}

Let $\mathcal{D}_{\text{cal}}$ contain factual claims with their scores. A
"trace" denotes a generation trajectory indexed by
$(\text{prompt\_id}, \text{seed})$. Our cal/test split protocol operates at
the trace level: every claim from a given trace is assigned to the same side
(stratified by category at the trace level; see Data Construction Pipeline
below). For each trace $i$ in category $c$, define a
\textbf{per-trace summary score}
\begin{equation}\label{eq:tracescore}
\bar{A}_i \;=\; \max_{\text{claim } k \in \text{trace } i,\ c_k = c}\, A_k,
\end{equation}
i.e., the maximum claim score over the trace's factual claims in category $c$
(or $-\infty$ if the trace has no factual claim in $c$). Let
$n_c^{\text{tr}}$ be the number of factual calibration traces in category
$c$. For each category, the trace-level per-category threshold is:
\begin{equation}\label{eq:Talphac}
\bar{T}_{\alpha, c} := \bigl\lceil (n_c^{\text{tr}} + 1)(1 - \alpha) \bigr\rceil\text{-th
order statistic of } \{\bar{A}_i : i \in \mathcal{D}_\text{cal},\ c_i = c,\ Y_i = \text{factual}\}.
\end{equation}

\begin{theorem}[Trace-level per-category conformal false alarm control]
\label{thm:layer2}
Assume $\{\bar{A}_1, \ldots, \bar{A}_{n_c^{\text{tr}}}, \bar{A}_{\text{new}}\}$
are the per-trace summary scores of factual traces in category $c$, and assume
\textbf{trace-level} exchangeability between calibration and test traces
within category $c$. Then
\[
\mathbb{P}\!\left(\bar{A}_{\text{new}} > \bar{T}_{\alpha, c}
   \mid Y_{\text{new}} = \text{factual},\ c\right)
   \;\le\; \alpha + \frac{1}{n_c^{\text{tr}} + 1}.
\]
\end{theorem}

\begin{proof}
Standard split conformal argument~\cite{vovk2005algorithmic}: under
exchangeability of the per-trace summary scores, the rank of
$\bar{A}_{\text{new}}$ among
$\{\bar{A}_1, \ldots, \bar{A}_{n_c^{\text{tr}}}, \bar{A}_{\text{new}}\}$ is
uniformly distributed over $\{1, \ldots, n_c^{\text{tr}}+1\}$. The probability
that it exceeds the
$\lceil (n_c^{\text{tr}}+1)(1-\alpha) \rceil$-th order statistic is at most
$\alpha + 1/(n_c^{\text{tr}}+1)$. The unit of exchangeability matches the
unit of the split (trace level), so the bound holds exactly in finite samples.
\end{proof}

\begin{remark}[Choice of unit: traces, not individual claims]
\label{rem:trace-unit-rationale}
We state Theorem~\ref{thm:layer2} at the trace level because that is the
unit of the cal/test split. Stating it at the individual-claim level
would assume claim-level exchangeability between calibration and test, but
multiple claims from the same trace share the token-level feature
trajectory, the cumulative wealth $M_t$ history, and the prompt; within-trace
dependence breaks claim-level exchangeability across the cal/test boundary.
The trace-level form delivers an exact finite-sample bound under the
assumption that actually holds in our protocol. The cost is reduced
resolution: each trace contributes one summary score rather than several
claim-level scores. We mitigate this by reporting trace-level conformal
calibration as the primary statistical guarantee, while also reporting
individual-claim AUROC and per-prompt miscoverage as informal
diagnostics (see Remark~\ref{prop:claim-level} and Group~B).
\end{remark}

\begin{remark}[Claim-level analogue threshold, informal]
\label{prop:claim-level}
For exploratory characterization, one may also define a per-claim
threshold $T_{\alpha, c}$ via the
$\lceil (m_c+1)(1-\alpha) \rceil$-th order statistic of the calibration
claim scores in category $c$ (with $m_c$ the number of factual calibration
claims). We do not state this as a Proposition or Theorem because the
claim-level exchangeability assumption it would require does not hold
under our trace-level split (within-trace dependence breaks individual-claim
exchangeability across cal/test). The analogous formal bound
$\mathbb{P}(A_{\text{new}} > T_{\alpha, c} \mid Y_{\text{new}} = \text{factual},\ c)
\le \alpha + 1/(m_c+1)$ therefore fails to apply directly; we use
$T_{\alpha,c}$ as an informal claim-level threshold whose deviation from
the formal bound is monitored empirically via per-prompt miscoverage with
Clopper--Pearson 95\%~confidence intervals (Group~B), flagging a
$(c, \text{prompt})$ cell only when its lower confidence bound exceeds
$\alpha + 1/(m_c+1)$ — this is a diagnostic flag, not a formal bound.
The claim-level scores are also the input to per-category AUROC and
PR-AUC characterizations (Layer~3); these characterizations are
reported as empirical correlation rather than as conformal-guaranteed
false-alarm control.
\end{remark}

\begin{remark}
The Layer~2 trace-level guarantee relies on exchangeability between
calibration and test \emph{traces} within each category---not on $\Hnull$
or any model-internal property. If exchangeability is broken (e.g.,
distribution shift in prompts or model version change), the guarantee
degrades. The exchangeability is empirical; we assess robustness via the
per-prompt miscoverage diagnostic in Group~B.
\end{remark}

\subsection{Data Construction Pipeline}

\subsubsection{Stage~1: Token-Level Feature Generation}

We use Qwen2.5-7B-Instruct as the target model. 25~prompts are divided into
five categories (C1--C5, see Table~\ref{tab:prompts}). For each prompt,
we generate 5~trajectories with random seeds
$\{42, 137, 256, 512, 1024\}$, yielding 125~generation trajectories.
Generation uses temperature $T = 0.7$, top-$p = 0.9$, and
$\texttt{max\_new\_tokens} = 1024$. Token features ($q_t$, $\|q_t\|_2^2$,
$p_t$, and auxiliary columns) are saved to a unified parquet file
(35,138~token rows). All downstream analyses read from this file;
the model is not re-invoked.

\subsubsection{Stage~2: Claim Extraction and Annotation}

\textbf{Phase A --- Claim extraction.} For each of the 125~generation
trajectories, we extract atomic factual claims with token-span mapping via the
DeepSeek V3.2~API. Token-to-character alignment is verified by round-trip
validation (100\% pass rate required). A pipeline fingerprint (SHA256 of
deterministic parameters) is embedded in each record. This produced
1,026~atomic claims.

\textbf{Phase B --- Multi-annotator annotation.} Each claim is independently
annotated by three LLM-based annotators:
DeepSeek V4 Pro (API),
GLM-4.7 (API), and
Claude Sonnet 4.6 (local CLI).
Each annotator labels every claim as one of:
\texttt{factual}, \texttt{hallucinated}, \texttt{uncertain}, or
\texttt{unverifiable}. Category-specific verification modules are used:
Knowledge Comparison (C1, C4), Recalculation (C2), Premise Detection (C3),
and Time-Anchored Fact Check (C5).

\textbf{Phase C --- Unanimity merge.} We apply a strict unanimity rule:
$n$:0 factual $\to$ factual (enters Layer~2 calibration set);
$n$:0 hallucinated $\to$ hallucinated (enters Layer~3 test set);
any dissent (factual vs.\ hallucinated) $\to$ needs human review;
any uncertain $\to$ excluded;
any unverifiable $\to$ excluded.
We report Cohen's $\kappa$ (pairwise), Fleiss' $\kappa$ ($n=3$), and the
label noise indicator $\eta_{\text{pair}}$. Claims are split into
calibration (70\%) and test (30\%) sets using per-category stratified
trace-level splitting (whole trajectory assigned entirely to calibration or
test, preventing information leakage across tokens within the same generation).

Final 3-annotator statistics from the completed annotation pipeline are
reported in \S\ref{ssec:annotation-stats}.

% ═══════════════════════════════════════════════════════════════════════
\section{Layer~3: Empirical Hallucination Bridge}

Layer~1 and Layer~2 are mathematically independent. The relationship between
token-level self-surprise and claim-level hallucination is an \textbf{empirical
question}---Layer~3 provides the quantitative answer.

\subsection{Evaluation Metrics}

For each claim with score $A$ and binary label
$Y \in \{0 = \text{factual}, 1 = \text{hallucinated}\}$, we compute:
\begin{itemize}
    \item \textbf{AUROC} and \textbf{PR-AUC} with 95\% bootstrap confidence
          intervals (30 resampling seeds);
    \item \textbf{Kolmogorov--Smirnov} $D$ statistic and $p$-value: tests
          whether the score distributions of factual vs.\ hallucinated claims
          differ;
    \item \textbf{Mann--Whitney} $U$ statistic and $p$-value: tests stochastic
          dominance of self-surprise scores for hallucinated claims.
\end{itemize}
All metrics are reported per category (C1--C5) and pooled. Bootstrap uses
trace-level resampling (sampling traces with replacement, not individual
claims) to preserve within-trace dependence structure.

\subsection{Expected Findings}

Based on the Channel~U / Channel~O structural analysis:
\begin{itemize}
    \item \textbf{C1, C2, C5 (Uncertainty-driven)}: AUROC substantially above
          0.5; self-surprise positively correlated with hallucination.
    \item \textbf{C3 (Implicit premise)}: AUROC near 0.5; the model absorbs
          false premises and outputs tokens with high $q_t(X_t)$, so
          self-surprise signals are weak.
    \item \textbf{C4 (Hedged hallucination)}: Possible \textit{phi-channel
          reversal}---self-surprise scores may be negatively correlated with
          hallucination in hedged expert domains.
\end{itemize}

The expected findings above are pre-registered predictions of the
Channel~U/O taxonomy; per-category numerical outputs are reported in
\S\ref{ssec:results-groupB} and \S\ref{ssec:results-groupC} (with full
numbers in the companion supplementary materials upon execution).

% ═══════════════════════════════════════════════════════════════════════
\section{Experimental Design}

\subsection{Prompt Categories}

\begin{table}[ht]
\centering
\caption{The 25~prompts across five categories
         (5~seeds each $=$ 125~trajectories).}
\label{tab:prompts}
\begin{tabular}{lll}
\hline
Category & Prompt IDs & Description \\
\hline
C1\_factual\_recall          & q001--q005 & Factual recall with time- \\
                             &             & \quad anchored verification \\
C2\_numerical\_reasoning      & q006--q010 & Math word problems \\
C3\_implicit\_premise          & q011--q015 & Questions embedding false \\
                             &             & \quad or tricky premises \\
C4\_domain\_expertise          & q016--q020 & Expert-level explanations \\
C5\_post-cutoff\_knowledge     & q021--q025 & Time-sensitive facts after \\
                             &             & \quad Qwen2.5's cutoff \\
\hline
\end{tabular}
\end{table}

\subsection{Model Configuration}
\begin{itemize}
    \item \textbf{Model}: Qwen/Qwen2.5-7B-Instruct (FP16)
    \item \textbf{Decoding}: $T = 0.7$, top-$p = 0.9$,
          $\texttt{max\_new\_tokens} = 1024$
    \item \textbf{Seeds}: $\{42, 137, 256, 512, 1024\}$
    \item \textbf{System prompt}: ``You are a helpful assistant.
          Always respond in English.''
    \item \textbf{Token suppression}: emoji, control characters, PUA, CJK,
          Hangul ($\sim$32K tokens suppressed; Lemma~\ref{lem:mean-zero}
          holds for any $\mathcal{F}_{t-1}$-measurable $q_t$, so suppression
          has no bearing on mean-zero).
\end{itemize}

\subsection{Group~A: Layer~1 Sanity Check}

\textbf{Objective}: Empirically verify Theorem~\ref{thm:layer1}.

\textbf{Method}: Run Algorithm~1 (see Appendix) on neutral-prompt generations
where $\Hnull$ holds by construction, computing
$\hat{\mathbb{P}}(\sup_t M_t \ge 1/\alpha)$ for
$\alpha \in \{0.01, 0.05, 0.1\}$, verifying that the empirical rejection rate
does not exceed $\alpha$.
\begin{itemize}
    \item \textbf{Positive control}: Inject a $\delta$-shift into $q_t$
          at a random token position and verify the e-process detects it.
    \item \textbf{Negative control}: Shuffle tokens randomly; verify the
          e-process does NOT trigger ($\Hnull$ preserved under permutation).
    \item \textbf{Per-category}: Verify $\Hnull$ holds uniformly across
          C1--C5, establishing that any cross-category variation in Layer~3
          is due to claim-level phenomena.
\end{itemize}

Empirical rejection rates and control results for Group~A are reported in
\S\ref{ssec:results-groupA} (pre-registered; numerical outputs in
companion supplementary materials upon execution).

\subsection{Group~B: Calibration and Detection Performance}

\textbf{Objective}: Validate Theorem~\ref{thm:layer2} and quantify
self-surprise--hallucination correlation per category.

\textbf{Method}: 30~random stratified trace-level splits (70/30). For each
split:
\begin{itemize}
    \item Compute the trace-level per-category threshold
          $\bar{T}_{\alpha, c}$ from cal-factual trace summary scores
          (Theorem~\ref{thm:layer2}, formal guarantee
          $\le \alpha + 1/(n_c^{\text{tr}}+1)$) and the informal
          claim-level threshold $T_{\alpha, c}$ from cal-factual claim
          scores (Remark~\ref{prop:claim-level}, informal reference
          $\alpha + 1/(m_c+1)$);
    \item Evaluate $\hat{\mathbb{P}}(\bar{A}_{\text{new}} >
          \bar{T}_{\alpha, c} \mid Y = \text{factual}, c)$ on test set
          (formal trace-level bound) and
          $\hat{\mathbb{P}}(A_{\text{new}} > T_{\alpha, c} \mid
          Y = \text{factual}, c)$ (informal claim-level reference);
    \item Compute AUROC, PR-AUC, KS, Mann--Whitney $U$ per category.
\end{itemize}
Bootstrap 95\% CIs across 30~splits.

\subsection{Group~C: Channel U / Channel O Demarcation}

\textbf{Objective}: Empirically characterize the boundary between the two
channels.

\textbf{Method}: Group categories by expected dominant channel:
Channel~U (C1, C2, C5) vs.\ Channel~O (C3, C4 hedged subset).
Compare AUROC between groups. A significant gap provides empirical support for
the theoretical demarcation. The C3/C4 deficit is not a failure of the
method---it is a \textit{measured characterization} of the structural
limitation.

\subsection{Group~D: Ablation Studies}

\textbf{Objective}: Quantify the contribution of each design choice.
\begin{itemize}
    \item \textbf{Score function}: Centered Brier residual vs.\ centered
          log-score $-\log q_t(X_t) - H(q_t)$ vs.\ v6~legacy score;
    \item \textbf{Sampling distribution}: $q_t$ (correct) vs.\ $p_t$ (raw
          softmax, to demonstrate that mis-specification breaks mean-zero);
    \item \textbf{Betting}: ONS vs.\ EMA plug-in vs.\ fixed $\beta = 0.1$;
    \item \textbf{Mixture weights}: Dyadic vs.\ geometric grid vs.\ single cohort;
    \item \textbf{Claim aggregator}:
          $\claimscore$ (Eq.~\ref{eq:aclaim}) vs.\
          $\claimscorelog$ (Eq.~\ref{eq:aclaillog}) vs.\
          global $\max_t M_t$ (incorrect baseline);
    \item \textbf{Conformal}: Pooled $T_\alpha$ vs.\ per-category $T_{\alpha,c}$.
\end{itemize}

Ablation results are reported in \S\ref{ssec:results-groupD} (pre-registered;
numerical outputs in companion supplementary materials upon execution).

% ═══════════════════════════════════════════════════════════════════════
\section{Results}
\label{sec:results}

This section reports two categories of results. \textbf{Stage~1 results}
(\S\ref{ssec:annotation-stats}) are complete: the annotation pipeline
(Phases A--C of the Data Construction Pipeline,
described above) has produced the labeled dataset on which all downstream
analyses operate, and we report per-annotator distributions,
inter-annotator agreement, the unanimity-merge yield, and per-category
calibration sample sizes.
\textbf{Stage~2 results} (\S\S\ref{ssec:results-groupA}--\ref{ssec:results-groupD})
are pre-registered analyses whose numerical outputs depend on the Phase~D
execution plan in \texttt{phase\_d/PIPELINE\_MODIFICATIONS\_PLAN\_v2.md}.
At submission time, the framework, the theory, the dataset, and the
pre-registered acceptance criteria (with their associated manifests and
thresholds) are in place; the numerical outputs of Groups~A--D and of the
streaming experiments (Path~B) are reported in the accompanying companion
supplementary materials upon execution. Three new honest disclosures
(\S\S\ref{ssec:llm-judge-circularity}--\ref{ssec:selection-on-agreement})
delimit the scope of these analyses regardless of their numerical outcome.

\subsection{Annotation Statistics}
\label{ssec:annotation-stats}

Of the 1{,}026~atomic claims extracted from 125~generation trajectories,
\textbf{667 claims} passed the 3-annotator unanimity merge rule (560 factual
and 107 hallucinated), with the remaining 359 excluded as needs-review for
one of the following reasons: any annotator chose \texttt{uncertain} or
\texttt{unverifiable}; the three annotators were split between
\texttt{factual} and \texttt{hallucinated}; or any annotator failed to
return a parseable response. The trace-level stratified split assigns
\textbf{447 claims to calibration} and \textbf{220 claims to test} across
82~calibration traces and 39~test traces.

Per-annotator label distributions (1{,}026 claims each) are: DeepSeek
V4~Pro 74.1\% factual / 13.2\% hallucinated / 4.4\% uncertain /
8.4\% unverifiable; GLM-4.7 76.3\% / 16.3\% / 1.9\% / 5.5\%; Claude
Sonnet~4.6 57.1\% / 16.2\% / 8.6\% / 18.1\%. Claude's higher use of
\texttt{uncertain} and \texttt{unverifiable} reflects a structurally
stricter rubric application; this is a known annotator-level behavior, not
a label-noise pathology.

Inter-annotator agreement: Fleiss' $\kappa$ ($n=3$, 4-label) is
\textbf{0.5763} (moderate). Pairwise Cohen's $\kappa$ (4-label) are
DeepSeek--GLM 0.70, DeepSeek--Claude 0.55, GLM--Claude 0.53. The binary
collapse (\texttt{factual} vs.\ other) gives DeepSeek--GLM
$\kappa_\text{bin} = 0.74$, indicating that the bulk of the 4-label
disagreement concerns the gradient between \texttt{uncertain} and
\texttt{factual}/\texttt{hallucinated} rather than the
factual/hallucinated dichotomy itself.

Per-category calibration sample sizes ($\alpha = 0.05$). The
\textbf{trace-level} cal counts $n_c^{\text{tr}}$ (number of factual
calibration traces per category, where a trace is factual-in-category $c$
if it has at least one factual claim in category $c$) determine the exact
Theorem~\ref{thm:layer2} bound $\bar{B}_c = \alpha + 1/(n_c^{\text{tr}}+1)$;
the \textbf{claim-level} cal counts $m_c$ (factual cal claims) determine
the informal Remark~\ref{prop:claim-level} reference value
$B_c = \alpha + 1/(m_c+1)$.

\begin{table}[ht]
\centering
\caption{Per-category claim-level calibration sample sizes $m_c$ (factual
calibration claims) and informal Remark~\ref{prop:claim-level}
reference values $B_c = \alpha + 1/(m_c+1)$ at $\alpha = 0.05$. Total calibration
claims per category (factual + hallucinated, in parentheses) are reported
alongside. The trace-level cal counts $n_c^{\text{tr}}$ used by
Theorem~\ref{thm:layer2} and their corresponding bounds
$\bar{B}_c = \alpha + 1/(n_c^{\text{tr}}+1)$ are computed during Phase~D
execution by \texttt{phase\_d/conformal\_layer2.py} from the trace-level
split audit and reported in the companion supplementary materials; based
on 82 total factual calibration traces stratified across 5~categories,
$n_c^{\text{tr}}$ will be substantially smaller than $m_c$, so the
trace-level bound $\bar{B}_c$ will be materially looser than the
claim-level bound $B_c$, an honest cost of matching the unit of
exchangeability to the unit of the cal/test split.}
\label{tab:cal-counts-bounds}
\begin{tabular}{lrrc}
\hline
Category & total cal claims & $m_c$ (factual) & $B_c$ \\
\hline
C1 Factual Recall              & 43  & 37  & 0.0763 \\
C2 Numerical Reasoning         & 86  & 86  & 0.0615 \\
C3 Implicit Premise            & 89  & 65  & 0.0652 \\
C4 Domain Expertise            & 170 & 156 & 0.0564 \\
C5 Post-Cutoff Knowledge       & 59  & 40  & 0.0744 \\
\hline
\end{tabular}
\end{table}

Total test claims per category: C1~22; C2~37; C3~44; C4~91; C5~26 (claim
level; trace-level test counts $\le$ claim counts per category and reported
in the companion supplementary materials upon execution). The structural
limitations on C1 (one hallucinated test claim) and C2 (zero hallucinated
test claims after trace-level split; the two unanimous-hallucinated C2
claims belong to a single trace placed entirely in calibration) are
discussed in \S\ref{ssec:c2-structural} and
\S\ref{ssec:effective-single-category}; the per-category test
factual/hallucinated breakdowns are derived from
\texttt{claim-level/claims\_final.jsonl} and reported in the companion
supplementary materials.

\subsection{Group~A: Theorem~\ref{thm:layer1} Verification (Pre-Registered)}
\label{ssec:results-groupA}

\textbf{Pre-registered analysis.} Under $\Hnull$, we evaluate the
empirical rejection rate $\hat{\mathbb{P}}(\sup_t M_t \ge 1/\alpha)$ for
$\alpha \in \{0.01, 0.05, 0.10\}$ across the 125~baseline trajectories.
The acceptance criterion (numerical tolerance accounting for finite-sample
fluctuation around the nominal $\alpha$ via a Clopper--Pearson upper CI)
is committed to \texttt{phase\_d/runs/<ts>/manifest.json} before the run,
per the pre-registration discipline in
\texttt{phase\_d/PIPELINE\_MODIFICATIONS\_PLAN\_v2.md} §6. \emph{Numerical
outputs are reported in the companion supplementary materials upon
execution.}
The Group~A positive control (a controlled $\delta$-shift injected at a
known token position) and the negative control (a fresh-seed regeneration
under identical decoding configuration) are similarly pre-registered. We
note that a shuffled-token control serves as a \emph{positive} control
rather than a negative control under our $q_t$-based score: permuting
tokens without re-running the model breaks
$X_t \mid \mathcal{F}_{t-1} \sim q_t$ (the permuted $X_t$ is from a
different filtration than the $q_t$ used to score it), so the e-process is
expected to trigger. The genuine negative control is the fresh-seed
regeneration, where $\Hnull$ holds by construction.

\subsection{Group~B: Detection Performance (Pre-Registered)}
\label{ssec:results-groupB}

\textbf{Pre-registered analysis.} We report per-category AUROC, PR-AUC,
Kolmogorov--Smirnov $D$, and Mann--Whitney $U$ on the test set, with 95\%
trace-level paired bootstrap confidence intervals (1{,}000 resamples,
seed=42, sampling at the trace level to preserve within-trace dependence).
The headline batch result is the pooled-by-percentile paired AUROC delta
of $\text{eprod}_{\beta=0.10}(s_t)$ versus $\text{mean}(\text{surprisal})$;
the pre-registered acceptance criterion is that this delta's 95\%
trace-level bootstrap CI excludes zero on the positive side and the
per-length-quartile signs are consistent across all length quartiles with
$\ge 5$ hallucinated test claims (degenerate quartiles excluded). Empirical
miscoverage diagnostics are computed at the
$(\text{category}, \text{prompt})$ cell level via Clopper--Pearson 95\%
intervals; a cell is flagged as exceeding the per-category bound
$\alpha + 1/(m_c+1)$ only when its \emph{lower} CI exceeds the bound
(strict, CI-based flag). \emph{Numerical results, including the 5-row
per-category table, the 6-baseline pooled headline, the length-stratified
breakdown, and the 25-cell prompt-conditional miscoverage report, are
reported in the companion supplementary materials upon execution per
\texttt{phase\_d/PIPELINE\_MODIFICATIONS\_PLAN\_v2.md} Angles~3 and~4.}
Power and structural caveats follow from
\S\ref{ssec:effective-single-category}: C1 and C2 are reported with
\texttt{power\_warning} flags, and the headline result is interpreted
predominantly through C5.

\subsection{Group~C: Channel Characterization (Pre-Registered)}
\label{ssec:results-groupC}

\textbf{Pre-registered analysis.} The Channel~U / Channel~O demarcation
(\S\ref{sec:channels}) generates a falsifiable per-category prediction:
the expected-Channel-U categories (C1, C5, and partially C2) should show
AUROC substantially above $0.5$, while the expected-Channel-O categories
(C3 and the false-premise subset of C4) should show AUROC near $0.5$ or
below. We report a Channel~U vs.\ Channel~O aggregate AUROC gap (with
trace-level paired bootstrap CI) and document any per-category deviation
from the channel prediction as a measured characterization of the
detectability boundary, not as a failure of the framework.
\emph{Numerical outputs are reported in the companion supplementary
materials upon execution.}

\subsection{Group~D: Ablation (Pre-Registered)}
\label{ssec:results-groupD}

\textbf{Pre-registered analysis.} The ablation grid covers: score function
(centered Brier residual vs.\ centered log-score vs.\ raw NLL); sampling
distribution ($q_t$ vs.\ $p_t$, where the $p_t$ variant is expected to
break Lemma~\ref{lem:mean-zero}); betting strategy (fixed $\beta=0.1$
vs.\ EMA plug-in; the ONS variant is deferred per
\texttt{phase\_d/PIPELINE\_MODIFICATIONS\_PLAN\_v2.md} §11); mixture
weights (dyadic vs.\ harmonic vs.\ single cohort); claim aggregator
($\claimscore$ vs.\ $\claimscorelog$); and conformal threshold (pooled
$T_\alpha$ vs.\ per-category $T_{\alpha,c}$). The primary betting/weight
combination for the main results is committed to
\texttt{phase\_d/runs/<ts>/manifest.json} before execution per the
pre-registration discipline in
\texttt{phase\_d/PIPELINE\_MODIFICATIONS\_PLAN\_v2.md} §6.4; all variants
are evaluated in the ablation. \emph{Numerical outputs are reported in the
companion supplementary materials upon execution.}

\subsection{Stage 4 Batch Headline: Gate 1 Outcome}
\label{ssec:results-batch}

We evaluate the multiplicative e-process candidate
$\text{eprod}_{\beta=0.10}(s_t)$ against six baseline aggregators
(mean-$s_t$, std-$s_t$, max-$s_t$, mean-surprisal, var-surprisal,
mean-entropy of $q_t$) on 220 test atomic claims (176 factual, 44
hallucinated). The headline pooled-by-percentile paired bootstrap delta
($\text{eprod}_{\beta=0.10}(s_t)$ versus $\text{mean(surprisal)}$) is
$\Delta_{\text{AUROC}} = -0.103$ with 95\% trace-level paired bootstrap CI
$[-0.246, +0.101]$ ($p_{\text{two-sided}} = 0.262$); all four non-degenerate
length quartiles show $\Delta_{\text{AUROC}} < 0$, consistent in direction
but not in statistical significance against zero.
Per-category AUROC: $C_1 = 0.381$ (underpowered, $n_h = 1$),
$C_2 = 0.407$ (underpowered, $n_h = 2$), $C_3 = 0.420$, $C_4 = 0.465$,
$C_5 = 0.536$ (only $C_5$ exceeds the random baseline of $0.5$ in point
estimate).
We conclude that the multiplicative e-process is \textbf{comparable to}
entropy/NLL baselines under our conformal calibration on this dataset; we
do not detect a length-uniform AUROC advantage at this sample size. The
conformal calibration framework (Theorem~\ref{thm:layer2}) and joint-
characterization framing (Theorem~\ref{thm:layer1} Remark) are preserved
regardless of this empirical outcome.

\subsection{Stage 7 Streaming Delay-Recall: Gate 2 Outcome}
\label{ssec:results-streaming}

Delay-recall comparison is not evaluated in this submission. The
test-hallucinated sample $n_h = 8$ trace-category records is structurally
insufficient for the paired-bootstrap power requirement at the
pre-registered binding cell $(\Delta=0.15, \rho=0.3)$; the Gaussian-copula
precheck at $n=40$ yields power $P = 0.176$ (Stage~6 precheck), and the
operational paired-Bernoulli test at $n_h = 8$ has structurally lower power
than this precheck.
A power-audit on the 8 traces (Stage~7 pre-audit) confirms: neither
Arm~A (Ville threshold $1/\alpha = 20$) nor calibrated CUSUM
($h = 5.21$, $k \approx 0$) fires alarms on any of the 8 traces — a
structural detectability finding consistent with the Channel-O-leaning
hallucinated subset (5/8 records are $C_5$ post-cutoff knowledge).
We report this as a framework-level structural property: the multiplicative
wealth-process Arm~A satisfies finite expected stopping time under any
alternative hypothesis with strictly positive expected score
$\mathbb{E}_{H_1}[s_t] > 0$ (corollary of Theorem~\ref{thm:layer1}'s
finite-time mean-exit property); empirical delay-recall comparison at this
dataset's $n_h = 8$ is left to future work with a Channel-U-dominant
benchmark and larger hallucinated sample.

\subsection{Stage 8 Peeking-FAR Three-Arm: Gate 3 Outcome}
\label{ssec:results-peeking}

We evaluate the three-arm peeking-FAR experiment on 31 test-factual traces
at $K^* = 25$ peeks per trace. Per-category calibration counts
$n_c^{\text{tr}} \in \{13, 14, 17, 17, 17\}$ fall below the floor
$\lceil 1/\alpha \rceil + 1 = 21$ required for non-vacuous
$(1-\alpha)$-quantile under per-category split conformal
(Theorem~\ref{thm:layer2}), placing the Path~B detector arms in a
vacuous-threshold regime ($c^* = c_A^* = +\infty$ for all five categories).

Condition~2 (Arm~A Ville bound, Theorem~\ref{thm:layer1}): empirical
$\hat{\text{FAR}}_A = 0/31$ does not violate the predicted bound
$\alpha = 0.05$; the 95\% Clopper--Pearson upper CI is $[0, 0.112]$,
consistent with true rates including the bound but also values up to
$0.112$ that would exceed it. We cannot certify the bound at 95\%
confidence at $n_{\text{test\_factual}} = 31$; this is finite-sample
non-identifiability, not bound consistency.

Condition~3 (Arm~A' per-category conformal bound): for all five categories,
Arm~A' empirical FAR is $0/n_{\text{test\_factual\_in\_c}}$ with
Clopper--Pearson 95\% upper CI exceeding the conformal-tightened bound
$\alpha + 1/(n_c^{\text{tr}}+1) \in [0.106, 0.121]$; the per-category
$\Delta^{\text{AV}}$ values are reported with this caveat. The bound miss
is a finite-sample CI-width artifact in the per-category
$n_{\text{test\_factual\_in\_c}} \in \{3, 5, 7, 8\}$ regime, not a
violation of the conformal calibration framework.

Condition~4 (Arm~B $K=1$ sanity): $\hat{\text{FAR}}_B(K=1) = 0/31$ with
Clopper--Pearson 95\% two-sided CI $[0, 0.112]$ contains $\alpha = 0.05$,
validating the $K=1$ calibration sanity check at this sample size.

Condition~5 (joint-effect gap $\Delta^{\text{AV}}$): the gap under the
(anytime-valid status, calibration-target) joint perturbation is not
detectable at $K^* = 25$ on this trace sample. $\Delta^{\text{AV}}$ point
estimate is $0.0000$ with paired bootstrap 95\% CI $[0, 0]$ (degenerate
due to vacuous thresholds on both Arms~A' and~B). Empirical content of
$\Delta^{\text{AV}}$ requires non-vacuous calibration thresholds, which
in turn require $n_c^{\text{tr}} \geq 21$ at $\alpha = 0.05$.

\subsection{Structural Limitation: C2 Zero-Hallucinated Problem}
\label{ssec:c2-structural}

Under unanimity merging, C2 (Numerical Reasoning) produces only 2
hallucinated claims in the entire dataset of 667 unanimous claims (the
remaining C2 disagreement-claims are excluded as \texttt{uncertain} or
\texttt{unverifiable}, as the annotators preferentially label numerical
errors as \texttt{uncertain} rather than \texttt{hallucinated}). After
trace-level split, both hallucinated C2 claims fall in the calibration
side, leaving \textbf{zero hallucinated test claims for C2}. This renders
per-category AUROC \emph{undefined} for C2 under the current annotation
protocol; we report C2's result as \texttt{undefined} with a
\texttt{power\_warning} flag rather than computing an uninformative
bootstrap-degenerate point estimate. A pre-registered sensitivity analysis
on C1 and C2 under a 2-of-3 majority merge rule (the C3, C4, C5 categories
keep unanimity) is reported in the appendix; see also
\S\ref{ssec:selection-on-agreement} for the implications of unanimity
selection on inferred AUROC.

\subsection{LLM-as-Judge Circularity in Annotation}
\label{ssec:llm-judge-circularity}

The ground-truth labels in this study are produced by three LLM
annotators (DeepSeek V4 Pro, GLM-4.7, Claude Sonnet 4.6), and the
score function $s_t = \|q_t\|_2^2 - q_t(X_t)$ is derived from a
fourth LLM (Qwen2.5-7B-Instruct). The headline experiment therefore
asks whether LLM-derived scores predict LLM-judged hallucination
labels. This protocol carries a partial-tautology risk: any pattern
shared between the score-model and the judge-models---e.g., both
finding the same tokens surprising due to overlapping pretraining
data---will inflate AUROC without measuring true hallucination
detection. We do not have a human gold standard to disambiguate this.
Per-category AUROC in this paper should be interpreted as a measure
of agreement between two LLM-derived signals on the surviving
unanimous-consensus claims, rather than as a measure of true
hallucination detection performance against an independent ground
truth. A human-annotated subset is left for future work.

\subsection{Effective Single-Category Empirical Result}
\label{ssec:effective-single-category}

The 5-category design (C1--C5) yields empirically interpretable
results for at most one category under the current annotation
protocol. C2 (Numerical Reasoning) has zero hallucinated test claims
after unanimity merging and trace-level split, making per-category
AUROC undefined. C1 (Factual Recall) has one hallucinated test
claim, making bootstrap confidence intervals structurally degenerate
(\texttt{power\_warning}=\texttt{underpowered}). C3 (Implicit Premise
Identification) is hypothesized by the Channel~U/O taxonomy to fall
predominantly into Channel~O (premise-acceptance hallucinations) and
may therefore yield AUROC near $0.5$. C4 (Domain-Specific Expertise)
is mixed Channel~U/O; the per-category result may show no advantage
or AUROC$<0.5$ on overconfident-recall cases. Only C5 (Post-Cutoff
Knowledge) is a relatively clean Channel~U scenario, but C5 carries
its own annotation circularity (the LLM annotators' knowledge cutoff
overlaps the time-anchored claims, so an annotator may label a claim
\texttt{factual} merely because it is consistent with the
annotator's own training-data knowledge, rather than verified against
an external source as of the anchor date~2026-05-13). The empirical
contribution of this paper is therefore narrower than the
5-category framing suggests: the C5 result is the cleanest test of
the framework's Channel~U coverage, and C1--C4 are reported as
supporting characterization with their respective acknowledged power
or structural limitations.

\subsection{Selection-on-Agreement Bias from Unanimity Merge}
\label{ssec:selection-on-agreement}

The unanimity merge rule excludes claims where any of the three
annotators chose \texttt{uncertain} or \texttt{unverifiable}, and
excludes claims with mixed \texttt{factual}/\texttt{hallucinated}
verdicts (Phase~2C of the data construction pipeline). This drops approximately
35\% of annotated claims (from 1{,}026 to 667 final records) and
retains a selected set on which all three LLM annotators reached
consensus---i.e., the cases that are systematically the
\emph{easiest} for LLMs to label confidently. Per-category AUROC on
this selected set therefore reflects detection performance on
unambiguous claims and systematically overestimates the AUROC that
would be obtained on the underlying claim population, which includes
the harder cases that triggered annotator disagreement. A 2-of-3
majority sensitivity analysis on C1 and C2 (appendix) partially
probes this effect but does not eliminate it, as the majority rule
relies on the same three LLM annotators. A more thorough probe would
require either human adjudication of the disagreed cases or a fourth
independent annotator from a different model family; both are left
for future work.

% ═══════════════════════════════════════════════════════════════════════
\section{Discussion}

\subsection{What the Framework Guarantees}
\begin{itemize}
    \item \textbf{Layer~1}: Under $\Hnull$, anytime-valid Type~I error control
          at the token level (Theorem~\ref{thm:layer1}).
    \item \textbf{Layer~2}: Under trace-level exchangeability,
          finite-sample false alarm rate control on factual
          \emph{traces} per category (Theorem~\ref{thm:layer2}).
    \item \textbf{Layer~3}: Quantitative empirical characterization of the
          self-surprise--hallucination correlation, per category.
\end{itemize}

\subsection{What It Does NOT Guarantee}
\begin{itemize}
    \item Layer~1 does not directly yield hallucination false alarm control
          ($\Hnull$ holds regardless of factuality).
    \item Layer~2 does not cover Channel~O (overconfident hallucinations).
    \item Layer~3 provides correlation evidence, not a mathematical guarantee
          for individual claim-level decisions.
    \item The claim-level analogue bound $\alpha + 1/(m_c+1)$
          (Remark~\ref{prop:claim-level}) does not hold formally under
          our trace-level cal/test split protocol; it is reported as an
          informal empirical reference monitored by per-prompt
          Clopper--Pearson miscoverage diagnostics, with cells flagged
          when their lower CI exceeds $\alpha + 1/(m_c+1)$ as a
          diagnostic rather than a formal-bound test. The formal Layer~2
          guarantee is the trace-level bound
          $\alpha + 1/(n_c^{\text{tr}}+1)$ in Theorem~\ref{thm:layer2}.
    \item The framework does not hold under arbitrary distribution shift
          (Layer~2 relies on exchangeability).
\end{itemize}

\subsection{Channel~O as Future Work}

Channel~O is the frontier. The present dataset characterization
(\S\ref{sec:channels}) places this evaluation in a Channel-O-leaning
regime: of 107~hallucinated atomic claims, the tight Channel-O proxy
($C_5$) accounts for 28\%, the wider proxy ($C_1 + C_5$) accounts for
34.6\%, and the widest proxy ($C_1 + C_2 + C_5$) accounts for 36.4\%.
The Stage~7 pre-audit on test-hallucinated traces (\S\ref{ssec:results-streaming})
provides direct corroboration: neither Arm~A nor calibrated CUSUM fires
on any of the 8 evaluated trace-category records, with 5/8 belonging to
$C_5$. The centered-Brier residual $s_t = \|q_t\|_2^2 - q_t(X_t)$ is
structurally insensitive to Channel~O failures — when the model assigns
high probability $q_t(X_t)$ to a wrongly-confident token, the residual is
near zero or negative, providing no detection signal regardless of stopping
rule. This is a documented limitation of the framework's scope on this
dataset, not a framework defect.

Promising future directions include:
(i)~triggering a Channel~O verifier (multi-sample agreement, retrieval lookup)
when Channel~U alarms;
(ii)~designing a claim-level confidence-vs-support mismatch detector;
(iii)~extending the conformal layer to a combined Channel~U~+~Channel~O score;
(iv)~constructing a Channel-U-dominant benchmark via prompts designed to
elicit uncertainty-driven hallucinations (numerical limits, multi-hop
factual recall, implicit-premise questions) to evaluate the framework
within its design scope.

\subsection{Relationship to Existing Methods}

SelfCheckGPT~\cite{manakul2023selfcheck} uses multi-sample consistency
post-hoc, addressing Channel~O but not real-time.
Semantic Entropy~\cite{farquhar2024detecting} is explicitly scoped to
confabulation; our framing is a subset.
FActScore~\cite{min2023factscore} relies on external knowledge, viewable as
a Channel~O retrieval-grounded implementation.
Our position is complementary: a real-time, anytime-valid,
internal-information detector for Channel~U, orthogonal to the above.

% ═══════════════════════════════════════════════════════════════════════
\section{Related Work}

\paragraph{LLM Hallucination Detection.}
Multi-sample consistency~\cite{manakul2023selfcheck}, semantic entropy for
confabulation~\cite{farquhar2024detecting}, atomic fact verification with
external knowledge~\cite{min2023factscore}, and internal state probe
methods~\cite{azaria2023internal} constitute the main threads. Our framework
complements them with real-time, assumption-lean, internal-information signals.

\paragraph{Strictly Proper Scoring Rules.}
The Bregman divergence characterization of proper scoring
rules~\cite{gneiting2007strictly} underpins our score function design.
Shao et al.~\cite{shao2024language} apply strictly proper scoring rules
to LLM token-level generation; our use of the centered Brier residual for
sequential testing is a different direction.

\paragraph{E-Processes and Anytime-Valid Inference.}
The game-theoretic statistics framework~\cite{ramdas2023gametheoretic,
howard2021timeuniform,waudbysmith2024estimating} provides the anytime-valid
machinery. The e-detector framework~\cite{shin2023edetectors} is the direct
source of our cohort~$+$~mixture construction; our use of it is an
instantiation with a specific score (centered Brier residual on the
post-temperature, post-top-$p$ sampling distribution) rather than a new
sequential-testing theorem, and we present Theorem~\ref{thm:layer1} as a
structural property of $M_t$ under $\Hnull$ rather than as a novel
theoretical contribution.

\paragraph{Closest prior work: Su et al.~\cite{su2026online}.}
Su, Wang, and Zhao apply e-processes to online LLM watermark detection.
The construction is the closest prior work to our Layer~1 in spirit, but
the differences are methodological rather than incremental: (i)~their null
hypothesis tests the presence/absence of a known watermarking scheme on
the token stream (an independence hypothesis on a pivotal statistic
depending on the watermark key), so their $H_0$ is falsifiable and the
e-process inflates under watermarked sequences; our null
$\Hnull: X_t \mid \mathcal{F}_{t-1} \sim q_t$ is a behavioral fact that
holds on every model trajectory by construction (Lemma~\ref{lem:mean-zero}),
so our e-process is a calibration object for a structural property of
$M_t$, not a falsifiable detector of an external attribute; (ii)~their
score is a pivotal statistic of the watermark hash function applied to
token outputs; ours is the centered Brier residual of $q_t$ on $X_t$,
which is meaningful whether or not any watermark is present and which is
linked through Theorem~\ref{thm:layer2} to a per-category conformal false
alarm rate; (iii)~Su et al.\ stop at the e-process layer (Layer~1 in our
nomenclature); we add a per-category split conformal calibration layer
(Theorem~\ref{thm:layer2}) and an empirical hallucination bridge layer
(Layer~3) on top, which together carry the hallucination-detection
performance claims of this work, with Layer~1 used to characterize
behavior under repeated peeking on factual traces rather than to deliver
detection. The complementary distinction is that watermark detection has
a non-trivial falsifiable null whereas hallucination detection (in our
formulation) does not; our framework therefore relies on the empirical
bridge (Layer~3) rather than on Layer~1 to do the detection work.

\paragraph{Conformal Prediction.}
Split conformal~\cite{vovk2005algorithmic,lei2018distribution} provides the
finite-sample guarantee in Layer~2. Our per-category stratified calibration
extends the standard split-conformal recipe.

% ═══════════════════════════════════════════════════════════════════════
\section{Conclusion}

We presented SHT, a three-layer framework for LLM hallucination detection:
(1)~token-level self-surprise e-process with anytime-valid Type~I error
control (Theorem~\ref{thm:layer1});
(2)~claim-level split conformal calibration with per-category false alarm
rate bounds (Theorem~\ref{thm:layer2}); and
(3)~empirical bridge quantifying the self-surprise--hallucination correlation
per category.
The three layers have distinct mathematical targets and do not collapse into
a single claim. We explicitly demarcate Channel~U (detectable by self-surprise)
from Channel~O (overconfident, structurally invisible to token-level scoring
rules) and declare the latter as explicit future work.

% ═══════════════════════════════════════════════════════════════════════
%  Appendix (unnumbered sections in LLNCS; use \section* or omit numbering
%  depending on page budget)
% ═══════════════════════════════════════════════════════════════════════
\section*{Appendix A: Algorithm~1 --- Token-Level Self-Surprise E-Process}

\begin{verbatim}
Input: {q_t}_{t>=1}, {X_t}_{t>=1}, weights {w_j}
       with sum_j w_j <= 1, B in (0,1), alpha.

 1: M <- sum_{j>=1} w_j          // includes inactive tail
 2: for t = 1, 2, ... do
 3:     // Step 1: choose bets (F_{t-1}-measurable)
 4:     for each active j <= t with w_j > 0 do
 5:         compute beta_t^(j) in [0,B]
 6:            using only s_1,...,s_{t-1}
 7:     // Step 2-3: observe token, compute score
 8:     observe X_t ~ q_t
 9:     s_t <- ||q_t||_2^2 - q_t(X_t)
10:     // Step 4: update cohort wealth
11:     for each active j <= t do
12:         Lambda^(j) <- Lambda^(j)*(1 + beta_t^(j)*s_t)
13:     // Step 5: activate cohort j = t+1
14:     M <- sum_{j<=t+1} w_j*Lambda^(j)
15:          + sum_{j>t+1} w_j
16:     if M >= 1/alpha then return alarm at tau = t
17: end for
\end{verbatim}

\section*{Appendix B: Pipeline Fingerprint Specification}

Every record in \texttt{claims\_final.jsonl} carries four SHA256 fingerprints
plus a schema version, providing end-to-end reproducibility binding from the
extraction config down to the underlying token-feature parquet.

\textit{Extraction fingerprint} (Layer~A): config hash of the claim extraction
run, covering prompt SHA, extractor model, extraction temperature, and planned
annotator set. Computed by \texttt{extract\_claims\_v2.compute\_fingerprint()}.

\textit{Merge fingerprint}: hash of (extraction fingerprint $+$ actual
annotator list $+$ merge rule version $+$ $n_{\text{annotators}}$).
This is an audit record of what merge produced; it is NOT a verification
target (the merge writes it; nothing upstream owns it).

\textit{Claims data fingerprint} (Layer~B): SHA256 of \texttt{claims\_v2.jsonl}
raw bytes. Detects post-hoc edits to \texttt{claim\_text},
\texttt{start\_token}, or \texttt{end\_token}.

\textit{Token features fingerprint} (Layer~C): SHA256 of
\texttt{token\_features.parquet} raw bytes. Detects silent parquet
regeneration with a different seed or configuration.

\textit{Fingerprint schema version}: tag identifying the hashing convention
(currently \texttt{3-layer-raw-bytes-v1}). Compared first; mismatch aborts
before any hash comparison, preventing cross-version equality bugs.

\textbf{Layer~3 verification contract.} Phase~D MUST, before any AUROC
computation, run \texttt{phase\_d/verify\_fingerprints.py} which performs
four checks in order: (i) schema-version equality (aborts on mismatch
before any hash comparison); (ii) Layer~A regression check --- recompute
via \texttt{extract\_claims\_v2.compute\_fingerprint()} and assert match
against the records' \texttt{extraction\_fingerprint} field; (iii) Layer~B
--- recompute claims-data fingerprint via
\texttt{claim-level/fingerprint\_utils.py} and assert match against
\texttt{claims\_data\_fingerprint}; (iv) Layer~C --- recompute
token-features fingerprint similarly and assert match against
\texttt{token\_features\_fingerprint}. Any mismatch $\to$ abort.

% ═══════════════════════════════════════════════════════════════════════
%  Bibliography (splncs04 style required by Springer LNCS)
% ═══════════════════════════════════════════════════════════════════════
\bibliographystyle{splncs04}
\bibliography{references}

\end{document}
